\label{sec:confidence-intervals}

\subsection{Sources of Uncertainty in Gingles Analysis}

Expert testimony in Section 2 litigation should quantify uncertainty, not
merely report point estimates.  There are three principal sources of
uncertainty in Gingles analysis:

\begin{enumerate}
  \item \textbf{Ecological inference uncertainty:} The WLS regression
    estimates precinct-level relationships from aggregate data.  Even a
    correctly specified model has estimation uncertainty from the finite
    sample of precincts.

  \item \textbf{Election selection uncertainty:} The choice of which elections
    to include in the analysis affects the estimates.  A robust methodology
    should demonstrate stability across the election sample.

  \item \textbf{District configuration uncertainty:} For Prong 1, the alignment
    score depends on the specific district configuration generated by the
    algorithm.  Different configurations satisfying population balance may
    yield different alignment scores.
\end{enumerate}

\subsection{Bootstrap Confidence Intervals from the Ensemble}

\citet{dellaLibera2026uncertainty} develops an ensemble methodology for
quantifying district configuration uncertainty.  The ensemble generates a
large number of algorithm runs with different random seeds, producing a
distribution of district configurations that satisfy the population-balance
and compactness constraints.  This distribution characterizes the space
of plausible maps.

For Prong 1, we report the 5th--95th percentile range of alignment scores
across the ensemble.  A Prong 1 finding is robust if the alignment score
exceeds 0.5 across at least 90\% of ensemble configurations.

For Prongs 2 and 3, we apply bootstrap resampling to the precinct-level
data~\citep{efron1979bootstrap}.  Each bootstrap resample draws precincts
with replacement and refits the WLS regression.  The 2.5th--97.5th percentile
range of the bootstrap distribution constitutes the 95\% confidence interval
for the regression coefficient.

\subsection{Reporting Standards for Expert Testimony}

We recommend the following reporting standards for Gingles expert testimony:

\begin{itemize}
  \item \textbf{Prong 1:} Report the alignment score $A_d$, minority CVAP
    percentage, and ensemble range.  If the ensemble range straddles the
    0.5 threshold, report the fraction of configurations satisfying Prong 1
    and the median configuration that provides the clearest visual
    illustration.

  \item \textbf{Prong 2:} Report $\hat{\beta}_{\text{min}}$ and 95\%
    bootstrap confidence interval for each election.  Report the range
    across elections and confirm that no single election drives the overall
    conclusion.  Minimum recommended sample: three contested elections in
    at least two election cycles.

  \item \textbf{Prong 3:} Report $\hat{\beta}_{\text{race}}$ with and without
    the partisan control variable, with HC3 standard errors and Holm-corrected
    $p$-values.  Report the LOO stability table.  Report the 95\% bootstrap
    confidence interval for the racial bloc-voting coefficient.  Minimum
    recommended sample: five general elections.
\end{itemize}

\subsection{Handling Ambiguous Cases}

Some minority communities will present ambiguous Gingles analysis: alignment
scores near 0.5, cohesion estimates near 0.7, or bloc-voting coefficients
that are positive but not statistically significant after Holm correction.
These cases require careful expert judgment about whether the evidence, taken
as a whole, supports a Prong finding.

Our general guidance is that ambiguous Prong 1 findings (alignment scores
0.45--0.55) should be supplemented with geographic narrative about the
continuity of the minority community and the reasons the algorithm places
its core in the proposed district.  Ambiguous Prong 2 or 3 findings should
be supplemented with additional elections from a longer time horizon,
noting whether the ambiguity reflects genuine instability in voting patterns
or merely limited data.
